\section{Inducing dimensions six and seven}
\label{sec:rank-six-seven-inducing}

Throughout this section $K$ is one of the eight non-Eisenstein
imaginary quadratic fields of class number one.  A source means
a connected finite-cover source with exact character field $K$,
actual Weil sixfold signature $(3,3)$, period rank nine, and
nonsplit actual polarization.  The inducing character field is
not assumed equal to $K$.

\begin{theorem}[D-BI-16: the remaining degree-six inducing group]
\label{D-BI-16}
If the coefficient representation of such a source is induced
from an irreducible representation $\tau$ of dimension six,
its faithful image and character field necessarily satisfy
\[
 K=\mathbf Q(i),\qquad
 L=\mathbf Q(\chi_\tau)=\mathbf Q(\zeta_{12}),\qquad
 R=G_{34},\qquad J=\mu_4G_{34}.                              \tag{BI.53}
\]
Here $R$ is the reflection normal closure in $J$; it has order
$39\,191\,040$ and center $\mu_6$, so $|J|=78\,382\,080$.
The actual intermediate base is rational, has $12\le r\le14$
genuine branch points, and has exactly two inertia entries outside
$R$.  Up to interchanging them, the complete necessary list of
moving-rank profiles is
\[
\begin{array}{c|c|c}
 r&\text{two outside moving ranks}&\text{inside moving ranks}\\\hline
14&(3,3)&1^{12}\\
13&(3,3)&1^{10},2\\
13&(3,4)&1^{11}\\
12&(3,3)&1^8,2,2\ \text{or}\ 1^9,3\\
12&(3,4)&1^9,2\\
12&(3,5)&1^{10}\\
12&(4,4)&1^{10}.
\end{array}                                                  \tag{BI.54}
\]
Exponents in this table denote repetitions, not powers of group
elements.  This table is necessary and does not claim that each
row occurs in a generating product-one tuple.

An outside rank-three inertia is an involution with spectrum
$1^3,(-1)^3$.  An outside rank-four inertia has spectrum
$1^2,(-1)^2,\alpha,\beta$, where, writing
$\zeta=\zeta_{12}$, the unordered residual pair is one of
\[
 \{i,i\},\ \{-i,-i\},\ \{i,-i\},\
 \{\zeta,\zeta^5\},\ \{\zeta^7,\zeta^{11}\},\
 \{\zeta,\zeta^{11}\},\ \{\zeta^5,\zeta^7\}.             \tag{BI.55}
\]
Its order is four or twelve, and its square is a rank-two
element of order two or six, respectively.  None of these
possibilities is asserted to occur merely from the eigenvalue
list.
\end{theorem}
\begin{proof}
Let $H$ be the inducing subgroup, so that the original
representation is $\operatorname{Ind}_H^G\tau$, and put
$L=\mathbf Q(\chi_\tau)$.  Induction gives $K\subset L$.
An imprimitive degree-six $\tau$ has an inducing block of
dimension one, two, or three.  Transitivity of induction and
D-MON, D-BI-01, and D-BI-05 exclude all three possibilities.
Thus $J$ is primitive.

On the actual intermediate base of genus $h$, the multiplicity
Euler equation is $D=18-12h$.  Genus one has at most six genuine
branches and at most six marked-base parameters; higher genus
is impossible.  Frobenius reciprocity for the actual variation
therefore forces
\[
          h=0,\qquad r\ge12,\qquad
          \sum_j\operatorname{rank}(1-g_j)=18.              \tag{BI.56}
\]
At least six of the inertias are reflections, since every other
entry has moving rank at least two.  The reflection normal
closure $R$ is consequently nontrivial.  Its fixed space is
$J$-stable and zero, and its irreducible root components would
give an imprimitivity system unless there were only one.  Thus
$R$ is irreducible.  If it were imprimitive, Cohen's
Proposition~2.13 would put its normalizer in a monomial group,
as none of its exceptional dimensions is six; this is already
excluded.  The real type $A_6$ is a scalar standard case excluded
by D-BI-08.  The complete reflection classification leaves
\[
                         R=W(E_6)\quad\text{or}\quad G_{34}.
\]
These are the published classification and corrected normalizer
inputs used in D-BI-13, with all lower-dimensional exceptions
already discarded.

First take $R=W(E_6)$.  Its normalizer is contained in
$\mathbf C^*R$: after the real-normalizer reduction, the
nontrivial diagram automorphism is $-w_0$, where $w_0$ is the
longest Weyl element.  Thus all matrices are scalar multiples
of the rational matrices of $R$.  Count the branch entries with
nonreal scalar phase by $t$.  Each moves at least three
dimensions, because a nonreal eigenvalue of a real matrix has
the same multiplicity as its distinct conjugate.  Therefore
$18\ge r+2t$, so $t\le3$.  There cannot be exactly one such
entry, by the scalar product-one relation modulo $\{1,-1\}$;
there cannot be none, since then $L$ would be real.

If $t=3$, the three moving ranks are all three.  Each scalar
phase has minimal polynomial over $\mathbf Q$ of degree at
most two, because its inverse has multiplicity three in a
degree-six rational characteristic polynomial.  Its order is
three, four, or six.  Modulo real signs the phases have order
three or two, and the product of three nonidentity such phases
can be one only if all three have order three.  Hence all
matrices are defined over $\mathbf Q(\omega)$, which cannot
contain the stipulated non-Eisenstein $K$.

If $t=2$, the two moving ranks sum to at most eight.  One of
the corresponding real matrices has the inverse scalar as an
eigenvalue of multiplicity at least two.  The scalar degree over
$\mathbf Q$ is at most three, and a nonreal root of unity of
such degree has degree two.  The two phases are inverse modulo
a real sign, and the other phases are real.  Since $R\subset J$,
the scalar matrices belong to $J$ and their traces show that
$L$ is exactly the quadratic scalar field.  Thus $L=K=\mathbf Q(i)$.
Multiplicity-one Frobenius reciprocity supplies a $K$-model of
$\tau$ via its central projector on the original $K$-model.
The polarized Shapiro comparison is then a positive rational
scalar comparison of the actual Hermitian forms.  The already
proved Gaussian $E_6$ source barrier forces splitness, a
contradiction.  This eliminates the real closure.

Now take $R=G_{34}$.  It has a model over $F=\mathbf Q(\omega)$,
all its reflections have order two, and its degrees are
$6,12,18,24,30,42$.  In particular its center is $\mu_6$ and
its character field is exactly $F$: the model gives one inclusion
and the trace of $\omega I$ gives the other.  These are Cohen's
root model/classification and Malle's model-field data, already
checked in D-BI-13.  The corrected normalizer theorem gives
$R\subset J\subset\mathbf C^*R$.  Write $J=CR$, where $C$
is the finite scalar subgroup of $J$.  Since every root of unity
in $F$ lies in $\mu_6\subset R$,
\[
             J\cap\operatorname{GL}_6(F)=R,
             \qquad J/R\simeq C/(C\cap R).                 \tag{BI.57}
\]

Let $t$ count inertias outside $R$.  There are at least two:
otherwise the cyclic quotient in (BI.57) is either trivial,
in which case $L=F$, or is not generated by a product-one tuple.
For an outside $g=\lambda w$, $w\in R$, one has
$\lambda\notin F$.  If its fixed dimension is $b$, the
minimal polynomial of $\lambda^{-1}$ over $F$ occurs at least
$b$ times in the degree-six characteristic polynomial of $w$.
In particular $b\le3$, so each outside inertia moves at least
three dimensions.  Again $t\le3$.

If $t=3$, each fixes three dimensions, and
$[F(\lambda_j):F]=2$.  The roots of unity giving degree two
over $F$ have order four or twelve, and all lie in
$\mathbf Q(\zeta_{12})\setminus F$.  Each therefore has the
same nonidentity class in $\mu_{12}/\mu_6$.  Three such
classes cannot have product one.  Hence $t=2$.

The two phases are inverse modulo $\mu_6$, all other entries
are in $R$, and generation gives
\[
                   J=\langle R,\lambda I\rangle,
                   \qquad L=F(\lambda).                    \tag{BI.58}
\]
The field equality follows from scalar traces and $F\subset L$.
The two moving ranks sum to at most eight, so one fixes at least
two dimensions and $[F(\lambda):F]\le3$.  The complete list
of scalar orders outside $F$ with this property is
\[
                        4,12,9,18.                         \tag{BI.59}
\]
Indeed $\varphi(d)\le6$, so one first lists
$1,2,3,4,5,6,7,8,9,10,12,14,18$.  Orders $1,2,3,6$ lie in $F$;
orders $5,8,10$ have degree four and orders $7,14$ degree six
over $F$, since their cyclotomic fields have trivial intersection
with $F$.  The remaining four
orders give $F(i)=\mathbf Q(\zeta_{12})$ or
$\mathbf Q(\zeta_9)$.  The latter field has the unique
quadratic subfield $F$, so it is excluded.  The only
non-Eisenstein imaginary quadratic subfield of $F(i)$ is
$\mathbf Q(i)$.  Since $\mu_6\subset R$, both orders four
and twelve generate the same group $\mu_4R$.  This proves
(BI.53).

Normalize every outside entry as $g=iw$ with $w\in R$.
The nontrivial automorphism of $F(i)/F$ pairs the eigenvalues
$i$ and $-i$ of $w$ with equal multiplicity.  Thus $g$ has
the eigenvalues $1$ and $-1$ with equal multiplicity $b$,
and its moving rank is $6-b$.  The two outside ranks are at
least three and have sum at most eight; together with (BI.56),
this gives exactly the table (BI.54) by partitioning the
remaining inside moving rank.  In particular $r\le14$.

For outside rank three, $b=3$ exhausts all six eigenvalues,
so $g$ is the stated involution.  For outside rank four,
$b=2$, and the characteristic polynomial of $w$ has the form
$(X^2+1)^2Q(X)$ with $Q$ monic of degree two over $F$.
Every root $\mu$ of $Q$ is a root of unity of degree at most
two over $F$.  It is not $i$ or $-i$, which would increase
$b$.  If it were primitive of order twelve, its conjugate
over $F$ would be $-\mu$, and
$\det w=-\mu^2\notin\{1,-1\}$.  This is impossible,
since $R$ is generated by order-two reflections.  The remaining
possibility is that both roots $\mu,\nu$ of $Q$ belong to
$\mu_6$, with $\mu\nu\in\{1,-1\}$.  Hence
\[
                  \alpha=i\mu,\qquad
                  \beta=i\varepsilon\mu^{-1},
                  \qquad\varepsilon\in\{1,-1\}.
\]
Listing the six choices of $\mu$ and two choices of
$\varepsilon$ gives precisely (BI.55).  The order and square
claims follow immediately from these exponents.  No unlisted
reflection power is used to exclude a possible rank-five entry.
\end{proof}

\paragraph{Signature and descent constraints.}
In (BI.54) the total Chevalley--Weil age must equal nine at every
embedding of $L$ extending either embedding of $K$.
The rank-three outside age is $3/2$; the rank-four outside
ages in (BI.55) are $3/2,2,5/2$.  Thus the $r=14$ row is
automatically balanced, the all-reflection-inside $(3,4)$ row
requires rank-four age two, the $(3,5)$ row requires outside
rank-five age $5/2$, and the $(4,4)$ row requires the two
outside ages to sum to four.  The other rows must still satisfy
the same total-age equation for their nonreflection inside entries.
There is one useful further restriction on the $r=12$, $(3,4)$
row.  Let $a$ be the age of its unique inside rank-two entry and
$b$ the age of its outside rank-four entry.  At the original
embedding the age equation is $a+b=3$.  The automorphism of $L$
fixing $i$ and conjugating $\omega$ changes inside eigenvalues
to their complex conjugates, so that $a$ becomes $2-a$.
For an outside matrix $g=iw$ it sends $g$ to $-\overline g$.
The seven pairs in (BI.55) show that this preserves the
rank-four age $b$; it also preserves the rank-three age $3/2$.
Balance at this second embedding therefore gives $2-a+b=3$.
Consequently $a=1$ and $b=2$ in this row as well.

The natural $L$-multiplicity abelian variety has dimension twelve.
Rational induction and $[L:K]=2$ require it, along the original
family, to be isogenous to the square of the original $K$-sixfold.
For $r=12$ the original rank-nine family dominates its intermediate
Hurwitz component, imposing a generic $M_2(K)$ endomorphism
condition there.  For $r=13$ or $14$ only a subvariety of
dimension at least nine with that condition is forced; dominance
of the entire intermediate component is not asserted.
Moreover D-BI-15 bounds the period rank by at most six whenever
the inducing subgroup is normal or its specified imprimitivity
system is unique.  All remaining cases must avoid those symmetry
hypotheses.

\begin{theorem}[D-BI-17: exclusion of degree-seven inducing blocks]
\label{D-BI-17}
No source in the stated non-Eisenstein class-number-one scope
has an irreducible inducing representation of dimension seven.
\end{theorem}
\begin{proof}
An imprimitive representation of prime dimension seven is
monomial, and transitivity of induction invokes D-MON.  Thus
the faithful inducing image $J$ is primitive.  The actual
intermediate Euler equation is $D=20-14h$, so the same genus
argument gives $h=0$, $r\ge12$, and $D=20$.
There are at least four reflection inertias.  The reflection
normal closure is irreducible by primitivity.  Cohen's normalizer
theorem excludes imprimitive reflection closures, and the real
$A_7$ case is excluded by D-BI-08.  The full reflection list
therefore leaves only $R=W(E_7)$.

Its normalizer is $\mathbf C^*R$, with rational Weyl matrices
and scalar center $\{1,-1\}$.  A nonreal scalar inertia fixes
at most three dimensions and moves at least four.  If $t$
counts these entries, $20\ge r+3t$ gives $t\le2$.
The nonreal character field and the scalar product-one relation
give $t\ge2$, hence $t=2$.  The two moving ranks sum to at
most ten, so one fixes at least two dimensions.  Its scalar
phase consequently has degree at most $\lfloor7/2\rfloor=3$
over $\mathbf Q$.  A nonreal root of unity of degree at most
three has degree two.  The reciprocal-phase relation, generation,
and scalar traces show that the entire character field $L$ is
this quadratic field.  Since $K\subset L$ and $K$ is
non-Eisenstein, $K=L=\mathbf Q(i)$.

As in the degree-six real case, multiplicity-one Frobenius
reciprocity supplies a genuine $K$-model of the inducing
representation, and polarized Shapiro identifies its actual
Hermitian form with the original one up to positive rational
scale.  The already proved Gaussian $E_7$ source barrier forces
splitness.  This contradiction excludes the last case.
\end{proof}
